\scanpage{14}
nous intéresserons surtout au cas (21), \unsure{où|ou} $\Delta=2$, sans nous interdire de spécialiser en car. 2 où on aura $\Delta=0$...

Revenons toujours à la situation $(E,\omega,\varphi)$, donnant naissance à $f:E\to E'$, $f':E'\to E$, d'où $\varphi':E'\times E'\longrightarrow\underline{\mathrm O}_S$, enfin à $[\, ,\,]_\varphi$ sur $E'$ et $[\, ,\,]_{\varphi'}$ sur $E'$. On va montrer que $f:E\to E'$ est un hom. d'algèbre, et \inserted{(si $\varphi$ sym)} on va faire opérer l'algèbre de Lie $E'$ sur $E$ par dérivations \inserted{d'algèbre} qui invarient $\varphi$, de façon que $E\xrightarrow{f}E'$ commute avec l'action de l'algèbre de Lie $E'$ sur $E$ et $E'$. Enfin, lorsque $\varphi$ \inserted{(ainsi que $\omega$)} est associé à une forme quadratique lisse spéciale, on va trouver une interprétation de la situation d'algèbres de Lie en termes d'un hom de schémas en groupes lisses $G\to G'$ et d'une opération de $G'$ sur $G$. \deleted{\unsure{Pensons|Prouvons|Prenons}}\inserted{Vérifions} d'abord
\[
 f[x,y]_{\varphi'}=
 \underset{\substack{\parallel\ \text{déf }(\unsure{12|11})\\ f\alpha'(\unsure{fx\wedge fy|f(x)\wedge f(y)})}}{[f(x),f(y)]_\varphi}\ ; \tag*{[49]}
\]
\deleted{en prenant le pr. \unsure{scalaires|scalaire} des deux mbs avec un $\unsure{Z|z}\in\Gamma(S,E)$, la formule devient, et posant $ {}^{t}f(z)=z'$, la formule à vérifier devient}

% archive sheet number: 13
